\begin{definition}
    Seien $a, b \in \Aa_{sa}$, dann schreiben wir
    $$ a \leq b \quad :\Longleftrightarrow \quad b - a \in \Aa_+. $$
\end{definition}

\begin{remark}{\ }
    \begin{enumerate}
        \item Die oben definierte Relation ist eine Halbordnung auf $\Aa_{sa}$. Reflexivität ist klar, gilt $a \leq b \leq c$, so folgt $b - a, c - b \in \Aa_+$ und damit $c - a \in \Aa_+$, also $a \leq c$. Für die Antisymmetrie sei $a \leq b$ und $b \leq a$. Dann ist $b - a \in \Aa_+$ und $-(b-a) \in \Aa_+$. Dann folgt
        $$ \sigma(b-a) \subseteq [0, +\infty), \quad -\sigma(b-a) = \sigma(-(b-a)) \subseteq [0, +\infty), $$
        also folgt $\sigma(b-a) = \{ 0 \}$ und damit $\Vert b - a \Vert = 0$.

        \item Gilt $a \leq b$ und ist $\lambda \geq 0$ so folgt $\lambda a \leq \lambda b$. Außerdem ist $-b \leq -a$.
        
        Ist $c \in \Aa_{sa}$, so folgt $a + c \leq b + c$.

        Ist zusätzlich $d \in \Aa_{sa}$ mit $c \leq d$, so folgt $a + c \leq b + d$.

        \item Sei $c \in \Aa_{sa}$, so gilt $c \leq \Vert c \Vert e$. Um dies einzusehen bemerken wir zunächst $\sigma(\Vert c \Vert e - c) \subseteq \R$. Ist $t>0$, so ist
        $$ -\Vert c \Vert > -\Vert c \Vert - t \notin \sigma(-c) = -\sigma(c), $$
        also erhalten wir
        $$ -t \notin \sigma(-c) + \Vert c \Vert = \sigma(\Vert c \Vert e - c), $$
        also $-t \in \rho(\Vert c \Vert e - c)$. Also ist $\sigma(\Vert c \Vert e - c) \subseteq [0, +\infty)$.

        \item Seien $0 \leq a \leq b$. Nach dem vorigen Punkt gilt auch $b \leq \Vert b \Vert e$. Also folgt
        $$ \Vert b \Vert - \sigma(a) = \sigma(\Vert b \Vert e - a) \subseteq [0, +\infty) $$
        und damit $\sigma(a) \subseteq (-\infty, \Vert b \Vert]$. Da $a$ positiv ist folgt $\sigma(a) \subseteq [0, \Vert b \Vert]$ und insgesamt $0 \leq \Vert a \Vert = r(a) \leq \Vert b \Vert$.

        \item Seien $a, b \in \Aa_+$, $\Vert a \Vert, \Vert b \Vert \leq 1$, also $0 \leq a \leq e$ und $0 \leq b \leq e$. Dann folgt $0 \leq e - a \leq e$ und damit $0 \leq e + b - a \leq 2e$. Also ist
        $$ 1 + \sigma(b-a) = \sigma(e + b - a) \subseteq [0, +\infty). $$
        Analog erhält man
        $$ 1 - \sigma(b-a) = \sigma(2e - (e+b-a)) \subseteq [0, +\infty) $$
        und damit $\sigma(b-a) \subseteq [-1, +1]$, also $\Vert a - b \Vert \leq 1$.

        \item Sei $x \in \Aa$, dann gilt $x^* x \in \Aa_+$. Um dies einzusehen, setze $b \coloneqq x^* x \in \Aa_{sa}$ und schreibe $b = b_+ - b_-$ mit $b_+, b_- \in \Aa_+$. Setze $y \coloneqq x b_-$, so ist
        $$ -y^* y = -b_- x^* x b_- = -b_- (b_+ - b_-) b_- = b_-^3. $$
        Damit folgt
        $$ \sigma(-y^* y) = \sigma(b_-^3) = \sigma(b_-)^3 \subseteq [0, +\infty). $$
        Schreibe nun $y = y_r + i y_i$, so ist $y^* = y_r - i y_i$. Damit berechnen wir
        $$ y^* y + y y^* = 2y_r^2 + 2y_i^2, $$
        und durch Umformen erhalten wir
        $$ y y^* = 2y_r^2 + 2y_i^2 - y^* y \geq 0. $$
        Beachte (vgl. Funktionalanalysis 2 Übung)
        $$ \sigma(y^* y) \setminus \{ 0 \} = \sigma(y^* y) \setminus \{ 0 \}, $$
        so ist die linke Seite in der linken Halbachse, die rechte Seite in der rechten Halbachse, also folgt $\sigma(y^* y) = \{ 0 \}$, damit $y^* y = 0$. Damit ist auch $b_-^3 = 0$ und damit (da $b_- \in \Aa_+$) auch $b_- = 0$, somit ist $x^* x \in \Aa_+$.

        Inbesondere erhalten wir die Darstellung
        $$ \Aa_+ = \{ x^* x : x \in \Aa \}. $$

        \item Sind $a, b \in \Aa_{sa}$, $y \in \Aa$, dann folgt aus $a \leq b$ das auch $y^* a y \leq y^* b y$. Um dies einzusehen schreiben wir
        $$ y^* b y - y^* a y = y^* (b-a) y = y^* p^2 y = (py)^* py $$
        für ein $p \geq 0$ und aus dem vorigen Punkt folgt die Behauptung.

        \item Sei $q \in \Aa_+$. Dann ist $q$ invertierbar genau dann wenn $0 \notin \sigma(q)$, was äquivalent (da $q \geq 0$) zu der Existenz eines $\delta > 0$ mit $\sigma(q) \subseteq [\delta, +\infty)$ ist, was wieder zu $\sigma(q - \delta e) \subseteq [0, +\infty)$ äquivalent ist, was $q \geq \delta e$ bedeutet.
        
        Ist $q \geq \delta e$ für $\delta > 0$, so folgt
        $$ (q^{-1})^* = (q^*)^{-1} = q^{-1}, $$
        womit $q^{-1}$ selbstadjungiert ist. Wegen der Darstellung für $\sigma(q^{-1})$ folgt auch $q^{-1} \geq 0$.

        \item Sei $q \in \Aa_+ \cap \Inv(\Aa)$ mit $q \geq \delta e$ für ein $\delta > 0$. Dann gilt
        $$ \sigma(\sqrt{q}) = \sqrt{\sigma(q)} \subseteq [\sqrt{\delta}, +\infty). $$
        Also ist auch $\sqrt{q} \in \Aa_+ \cap \Inv(\Aa)$. Wegen
        $$ (\sqrt{q}^{-1})^2 = q^{-1} $$
        gilt auch $\sqrt{q^{-1}} = (\sqrt{q})^{-1}$.

        \item Sei $0 \leq a \leq b$, wobei $a \in \Inv(\Aa)$. Zunächst gibt es ein $\delta > 0$ mit $\delta e \leq a \leq b$, womit auch $b \in \Inv(\Aa)$ folgt. Dann gilt
        $$ 0 \leq (\sqrt{a})^{-1} a (\sqrt{a})^{-1} \leq (\sqrt{a})^{-1} b (\sqrt{a})^{-1} \eqqcolon p. $$
        Insbesondere gilt $0 \leq (p-e)$ und damit
        $$ 0 \leq (\sqrt{p})^{-1} (p-e) (\sqrt{p})^{-1} = e-p^{-1}. $$
        Damit folgt $\sqrt{a} b^{-1} \sqrt{a} \leq e$ und damit nach Multiplizieren mit $(\sqrt{a})^{-1}$
        $$ 0 \leq b^{-1} \leq a^{-1}. $$

        \item Seien $0 \leq a \leq b$, dann folgt\footnote{Im Allgemeinen folgt \textbf{nicht} $0 \leq a^2 \leq b^2$!} $0 \leq \sqrt{a} \leq \sqrt{b}$. Um das zu zeigen setze $p \coloneqq \sqrt{a}, q \coloneqq \sqrt{b}$, so reicht es $q - p \geq 0$ zu zeigen. Äquivalent dazu müssen wir für beliebiges $t > 0$ zeigen
        $$ te + q - p \in \Inv(\Aa). $$
        Dazu setzen wir $x_t \coloneqq (te + q + p)(te + q - p)$. Dann ist $x_t^* = (te + q - p)(te + q + p)$. Schreibe nun $x_t = c_t + i d_t$ für gewisse $c_t, d_t \in \Aa_{sa}$, dann ist
        $$ 2 c_t = x_t + x_t^* = 2 t^2 + 4 t q + 2(b-a) \geq 2 t^2 e. $$
        Insbesondere ist $c_t$ positiv und invertierbar. Nach dem Spektralabbildungssatz folgt
        $$ \sigma(e + i \sqrt{c_t}^{-1} d_t \sqrt{e_t}^{-1}) = 1 + i \sigma(\sqrt{c_t}^{-1} d_t \sqrt{c_t}^{-1}) \subseteq 1 + i \R. $$
        Also ist das linksstehende Element, welches gleich $ \sqrt{c_t}^{-1} x_t \sqrt{c_t}^{-1} $ ist, invertierbar. Inbesondere sind auch $x_t$ und damit $x_t^*$ invertierbar. Nun ist
        $$ x_t^{-1} (te + q + p)(te + q - p) = e $$
        und
        $$ (te + q - p)(te + q + p)(x_t^*)^{-1} = e, $$
        womit $te + q - p$ invertierbar ist.
    \end{enumerate}
\end{remark}